\documentclass[11pt,a4paper]{article}

\def\courseTitle{Industrial Security (Bachelor)}
\def\sectionNumber{04}
\def\sectionTitle{Industrial Protocols}
\def\sheetKind{Solution Sheet}

% =====================================================================
% Shared preamble for exercise sheets and solution sheets.
%
% Define BEFORE \input{}-ing this file:
%   \def\courseTitle{...}
%   \def\sectionNumber{...}
%   \def\sectionTitle{...}
%   \def\sheetKind{Exercise Sheet}   or  Solution Sheet
%
% Optional:
%   \def\courseAuthor{...}
%
% The caller must issue \documentclass{...} (typically article) BEFORE
% loading this preamble.
% =====================================================================

\usepackage[utf8]{inputenc}
\usepackage[T1]{fontenc}
\usepackage[english]{babel}
\usepackage[a4paper, margin=2.2cm]{geometry}
\usepackage{graphicx}
\usepackage{listings}
\usepackage{xcolor}
\usepackage{tikz}
\usepackage{booktabs}
\usepackage{array}
\usepackage{enumitem}
\usepackage{titlesec}
\usepackage{fancyhdr}
\usepackage{hyperref}
\usepackage{amsmath}
\usepackage{amssymb}
\usepackage{tcolorbox}
\tcbuselibrary{skins, breakable}
\usepackage{needspace}

% Discourage solo lines split across pages.
\widowpenalty=10000
\clubpenalty=10000
\raggedbottom

% =====================================================================
% Shared TikZ styles for the Industrial Security lecture series.
% Loaded by every slide deck and exercise sheet that uses TikZ.
% =====================================================================

\usetikzlibrary{
  arrows.meta,
  shapes,
  shapes.geometric,
  shapes.symbols,
  shapes.multipart,
  positioning,
  fit,
  calc,
  backgrounds,
  decorations.pathmorphing,
  decorations.pathreplacing,
  decorations.text,
  patterns,
  matrix,
  shadows
}

% --- colour palette --------------------------------------------------
\definecolor{isOT}{HTML}{1F4E79}     % OT (deep blue)
\definecolor{isIT}{HTML}{6F2DA8}     % IT (purple)
\definecolor{isSafety}{HTML}{C00000} % safety red
\definecolor{isNet}{HTML}{2E7D32}    % network green
\definecolor{isAttack}{HTML}{B23A48} % attack red
\definecolor{isAccent}{HTML}{E08E0B} % amber
\definecolor{isMute}{HTML}{6B6B6B}   % grey
\definecolor{isLight}{HTML}{EDF2F8}  % light background
\definecolor{isPLC}{HTML}{2C5282}
\definecolor{isHMI}{HTML}{2F855A}
\definecolor{isSCADA}{HTML}{6B46C1}

% --- generic node styles --------------------------------------------
\tikzset{
  isnode/.style={
    draw, rounded corners=2pt, minimum height=8mm, minimum width=18mm,
    font=\small, align=center, thick
  },
  isbox/.style={isnode, fill=isLight},
  plc/.style={isnode, fill=isPLC!15, draw=isPLC, thick},
  hmi/.style={isnode, fill=isHMI!15, draw=isHMI, thick},
  scada/.style={isnode, fill=isSCADA!15, draw=isSCADA, thick},
  sensor/.style={isnode, fill=isNet!10, draw=isNet, minimum height=6mm, minimum width=14mm, font=\scriptsize},
  actuator/.style={isnode, fill=isAccent!15, draw=isAccent, minimum height=6mm, minimum width=14mm, font=\scriptsize},
  firewall/.style={isnode, fill=isAttack!15, draw=isAttack, thick},
  server/.style={isnode, fill=isMute!15, draw=isMute!80, thick},
  switch/.style={isnode, fill=isLight, draw=black!60, minimum height=6mm, minimum width=14mm, font=\scriptsize},
  workstation/.style={isnode, fill=white, draw=isOT, thick},
  attacker/.style={isnode, fill=isAttack!20, draw=isAttack, thick, font=\small\bfseries},
  inetnode/.style={draw, ellipse, minimum width=24mm, minimum height=10mm, font=\scriptsize, fill=isLight, thick},
  process/.style={isnode, fill=isOT!15, draw=isOT, thick, minimum width=24mm},
  zone/.style={
    draw, rounded corners=4pt, thick, dashed, inner sep=4mm
  },
  conduit/.style={-{Stealth[length=2.5mm]}, thick, isNet!80!black},
  attack/.style={-{Stealth[length=2.5mm]}, thick, isAttack, dashed},
  flow/.style={-{Stealth[length=2.5mm]}, thick, isOT!70!black},
  netline/.style={thick, black!60},
  axisline/.style={-{Stealth[length=2mm]}, thick, black!70},
  tag/.style={font=\scriptsize\itshape, fill=white, inner sep=1pt},
  level/.style={
    draw, fill=isLight, minimum width=0.85\linewidth, minimum height=8mm,
    font=\small, anchor=west
  },
  brace/.style={decorate, decoration={brace, amplitude=4pt, raise=2pt}},
  legendbox/.style={draw, rounded corners=2pt, fill=isLight, inner sep=2mm, font=\scriptsize, align=left}
}

% --- handy macros ----------------------------------------------------
% \plcicon, \hmiicon etc as simple labelled boxes; useful inline.
\newcommand{\plcicon}[1][PLC]{\tikz[baseline=-0.6ex]{\node[plc, minimum width=10mm, minimum height=5mm, font=\scriptsize, inner sep=1pt]{#1};}}
\newcommand{\hmiicon}[1][HMI]{\tikz[baseline=-0.6ex]{\node[hmi, minimum width=10mm, minimum height=5mm, font=\scriptsize, inner sep=1pt]{#1};}}
\newcommand{\fwicon}[1][FW]{\tikz[baseline=-0.6ex]{\node[firewall, minimum width=10mm, minimum height=5mm, font=\scriptsize, inner sep=1pt]{#1};}}


\providecommand{\courseAuthor}{Industrial Security Lecture Series}
\providecommand{\courseInstitute}{Department of Computer Science}
\providecommand{\companyLogo}{logo}

% --- Corporate branding overrides ------------------------------------
% Picks up logo / author / brand colours from the repo-root branding.tex
% when present; falls back to the defaults above otherwise.
\IfFileExists{../../../branding.tex}{% =====================================================================
% Corporate / institutional branding -- single file to customise the
% look of the entire course (slides, notes, exercises, labs).
%
% HOW TO REBRAND IN UNDER A MINUTE:
%   1. Replace `branding/logo.pdf` with your own logo
%      (PDF preferred; PNG and JPG also work -- name it `logo.<ext>`).
%   2. (Optional) change the two colours below to your brand palette.
%   3. (Optional) override the author/institute lines below.
%   4. Re-run `make` at the repo root. That's it.
%
% Everything in this file has sensible defaults; you can delete a line
% to fall back to the built-in defaults, or comment it out.
%
% This file is loaded automatically by the slides, exercise, and lab
% preambles -- you never need to \input it manually.
% =====================================================================

% --- Logo --------------------------------------------------------------
% Path is resolved from each leaf-build directory; the default points at
% the shared `branding/` folder at the repo root. If you put your logo
% somewhere else, set the full or relative path here (without extension):
\renewcommand{\companyLogo}{../../../branding/logo}

% --- Author / institute (appears on the title page footer) ------------
\renewcommand{\courseAuthor}{}
\renewcommand{\courseInstitute}{}

% --- Brand colours ----------------------------------------------------
% slideAccent      = primary brand colour (title bands, accent rule)
% slideSecondary   = secondary brand colour (section badge, highlight)
% Both use HTML hex codes. Keep enough contrast against white text.
\definecolor{slideAccent}{HTML}{00205B}      % deep navy
\definecolor{slideSecondary}{HTML}{00A0DC}   % light blue accent
}{}

% --- listings -------------------------------------------------------
\definecolor{codebg}{rgb}{0.96,0.96,0.96}
\lstset{
  backgroundcolor=\color{codebg},
  basicstyle=\ttfamily\footnotesize,
  breaklines=true,
  frame=single,
  numbers=left,
  numberstyle=\tiny\color{black!50},
  showstringspaces=false,
  tabsize=2
}

% --- header / footer ------------------------------------------------
\pagestyle{fancy}
\fancyhf{}
\renewcommand{\headrulewidth}{0.4pt}
\renewcommand{\footrulewidth}{0pt}
\lhead{\small \courseTitle}
\rhead{\small Section \sectionNumber{} -- \sheetKind}
\cfoot{\thepage}

% --- task / solution boxes -----------------------------------------
% `before={...}` guards every task/solution box with \needspace so the
% box doesn't start with only one or two lines of breathing room at the
% bottom of a page. If less than ~9 lines remain, LaTeX bumps the box
% to the next page; otherwise it starts here and breaks normally if it
% runs long.
% Boxes are `breakable` so very long solutions don't blow the page,
% but `before={\needspace{14\baselineskip}}` pushes them to a fresh
% page whenever fewer than ~14 lines remain. That covers the vast
% majority of single-question splits in practice.
\newtcolorbox{taskbox}[1]{
  enhanced, breakable,
  colback=isLight, colframe=isOT,
  fonttitle=\bfseries\color{white},
  coltitle=white,
  colbacktitle=isOT,
  title={#1},
  arc=2pt, boxrule=0.6pt,
  left=2mm, right=2mm, top=2mm, bottom=2mm,
  before={\par\vspace{2mm}\needspace{14\baselineskip}\par},
  after={\par\vspace{2mm}}
}

\newtcolorbox{solutionbox}{
  enhanced, breakable,
  colback=isNet!5, colframe=isNet,
  fonttitle=\bfseries\color{white},
  coltitle=white,
  colbacktitle=isNet,
  title={Solution},
  arc=2pt, boxrule=0.6pt,
  left=2mm, right=2mm, top=2mm, bottom=2mm,
  before={\par\vspace{2mm}\needspace{14\baselineskip}\par},
  after={\par\vspace{2mm}}
}

\newtcolorbox{hintbox}{
  enhanced, breakable,
  colback=isAccent!10, colframe=isAccent,
  fonttitle=\bfseries\color{white}, coltitle=white, colbacktitle=isAccent,
  title={Hint},
  arc=2pt, boxrule=0.6pt
}

\newtcolorbox{warnbox}{
  enhanced, breakable,
  colback=isAttack!8, colframe=isAttack,
  fonttitle=\bfseries\color{white}, coltitle=white, colbacktitle=isAttack,
  title={Caution},
  arc=2pt, boxrule=0.6pt
}

% --- section formatting --------------------------------------------
\titleformat{\section}{\Large\bfseries\color{isOT}}{\thesection}{0.5em}{}
\titleformat{\subsection}{\large\bfseries\color{isOT!80!black}}{\thesubsection}{0.5em}{}

% --- title block macro ----------------------------------------------
\newcommand{\sheetheader}{%
  \begin{center}
    {\Large\bfseries \courseTitle}\\[0.3em]
    {\large \sheetKind{} -- Section \sectionNumber}\\[0.2em]
    {\large \sectionTitle}\\[0.2em]
    \rule{\linewidth}{0.4pt}
  \end{center}
  \vspace{0.5em}
}


\begin{document}
\sheetheader

\noindent\textbf{Note to graders.}
Each model answer is intentionally short. Award full credit for a different
phrasing only if the student names the same wire-level mechanism or the same
compensating control layer; ``apply security'' or ``use a firewall'' without
specifics is not full credit.

\vspace{0.6em}

\begin{solutionbox}
\textbf{Exercise 4.1 -- Decode a Modbus TCP Request.}
MBAP header: Transaction ID \texttt{0x0001} (client correlation cookie),
Protocol ID \texttt{0x0000} (always zero for Modbus), Length \texttt{0x0006}
(number of bytes from the Unit ID onward), Unit ID \texttt{0x11} (decimal
\texttt{17}, the slave or gateway target).
Function code \texttt{0x03} = \emph{Read Holding Registers}, which returns
16-bit values from the holding-register address space.
Starting register \texttt{0x006B} (decimal \texttt{107}), quantity
\texttt{0x0003}, so registers \texttt{107..109} are read.
A successful response carries MBAP (7\,B) $+$ FC (1\,B) $+$ byte count
(1\,B) $+$ $3\!\times\!2$\,B of register data $=$ 15 bytes total; the MBAP
\texttt{Length} field on the response is \texttt{0x0009} (unit + FC + byte
count + 6 data bytes).
\textit{Rubric: $-1$ for each field misnamed; $-2$ if Length is reported as
12 instead of 9.}
\end{solutionbox}

\begin{solutionbox}
\textbf{Exercise 4.2 -- EtherNet/IP Scan Took the Line Down.}
(1) The CIP stack on an older ControlLogix CPU is single-threaded for
explicit messaging; a burst of \texttt{ListIdentity}/\texttt{GetAttributesAll}
requests can starve the scan task. In the PCAP, look for tens of explicit
TCP/\texttt{44818} sessions opened in rapid succession from the scanner IP
with no implicit UDP/\texttt{2222} traffic interrupted.
(2) Connection-table exhaustion: each scan probe opens a CIP connection;
ControlLogix has a finite Class 3 connection pool. Look for repeated
\texttt{ForwardOpen}/\texttt{ForwardClose} pairs and, eventually,
\texttt{ForwardOpen} responses with status \texttt{0x01}
(\emph{Connection Failure}).
(3) A malformed or borderline-CIP request triggered an assertion in the
firmware (this is the historical Project Basecamp / OT:ICEFALL pattern).
Look for the last request before the silence; if the same byte pattern
reproduces the freeze on a lab device, this is the cause.
\textit{Rubric: any three plausible OT-specific causes are acceptable; deduct
if the student blames generic ``DoS'' without naming a CIP-side resource.}
\end{solutionbox}

\begin{solutionbox}
\textbf{Exercise 4.3 -- Hardening a Mosquitto Broker.}
\begin{verbatim}
listener 8883
cafile   /etc/mosquitto/ca.crt
certfile /etc/mosquitto/server.crt
keyfile  /etc/mosquitto/server.key
require_certificate true
allow_anonymous false
password_file /etc/mosquitto/passwd
\end{verbatim}
Any four of the lines above count. \texttt{listener 8883} moves the broker
off plaintext \texttt{1883/tcp}; the three TLS lines enable confidentiality
and server authentication; \texttt{require\_certificate true} forces
\emph{client} authentication (mutual TLS) so an attacker cannot just publish
with a stolen password; \texttt{allow\_anonymous false} together with
\texttt{password\_file} blocks the anonymous default that Shodan finds on
hundreds of thousands of brokers.
\textit{Rubric: full credit only if the student justifies why each line
matters; ``it is more secure'' is not a justification.}
\end{solutionbox}

\begin{solutionbox}
\textbf{Exercise 4.4 -- ``Use OPC UA, It Is Secure''.}
The four configuration choices are:
(1) \emph{SecurityPolicy}: must be \texttt{Basic256Sha256} or
\texttt{Aes256\_Sha256\_RsaPss}. The insecure default is
\texttt{SecurityPolicy = None}, accepted by many configurators and visible
on the wire as un-signed UA messages.
(2) \emph{MessageSecurityMode}: must be \texttt{SignAndEncrypt}. The
insecure default is \texttt{None} or \texttt{Sign}-only.
(3) \emph{User authentication}: must use username/password or X.509 user
tokens with a real identity. The insecure default is the
\texttt{Anonymous} user token.
(4) \emph{Certificate trust list}: the server's trust list must contain only
expected client certificates, and self-signed certificates must \emph{not}
be auto-accepted. The insecure default is ``trust on first use'' with
auto-acceptance enabled for commissioning.
\textit{Rubric: accept ``application-instance certificate lifecycle / GDS''
as a substitute for (4); deduct if the student lists OPC UA features
(e.g.\ ``namespaces'') instead of security choices.}
\end{solutionbox}

\begin{solutionbox}
\textbf{Exercise 4.5 -- Why Modbus Has No Authentication.}
Modbus was designed in 1979 for short, point-to-point RS-485 loops inside
one cabinet. The physical layer \emph{was} the security boundary: if you
were on the wire you were trusted, and CPU cycles on the 8-bit
microcontrollers of the day were too scarce for any kind of crypto. Adding
authentication later would have broken every installed device, so the
specification was frozen.
The modern compensating-control pattern is a layered one: (a) network
segmentation that places Modbus devices in their own zone, (b) a stateful or
deep-packet-inspection firewall at the conduit that allowlists function
codes per source IP, (c) optionally Modbus Security (TCP/802 with TLS,
specified in 2018) where both endpoints support it, and (d) monitoring on
the segment for unexpected write function codes
(\texttt{0x05/0x06/0x0F/0x10}).
\textit{Rubric: full credit needs both the ``trust boundary was the wire''
historical argument and at least two compensating layers.}
\end{solutionbox}

\begin{solutionbox}
\textbf{Exercise 4.6 -- Reading a DNP3 Trace.}
Operationally, outstation \texttt{5} is pushing event data to master
\texttt{1} without being polled. This is the canonical DNP3 unsolicited
response pattern used so that, e.g., a relay can report a breaker trip
immediately instead of waiting for the next master cycle. Six class-1 events
in 50\,ms is plausible for a real fault (a single bay can generate a burst
of state changes).
The anomaly worth checking is whether the burst correlates with anything in
the SCADA event log: if the master reports no operator action and no
upstream fault, the unsolicited burst could be a spoofed or replayed
sequence (DNP3 without SAv5 has no authentication, so any host on the
segment can forge a frame with source address \texttt{5}). Confirm by
checking the application-layer sequence numbers for gaps or duplicates, by
correlating with the IED's own SOE log, and by checking whether SAv5
challenge/response is configured on the link.
\textit{Rubric: accept any anomaly that ties to spoofing/replay or to
event-storm overload; deduct if the student treats unsolicited responses
themselves as the anomaly.}
\end{solutionbox}

\begin{solutionbox}
\textbf{Exercise 4.7 -- Profinet vs.\ Modbus: Fail-Closed?}
Profinet fails closed first. Profinet RT runs a cyclic exchange with a
configured update time (here 4\,ms) and a watchdog factor (typically 3,
giving a $\sim$12\,ms timeout); an 800\,ms outage is two orders of magnitude
longer than the watchdog, so the controller drops the connection, the
device enters its configured fail-safe substitute values, and the
controller logs a diagnostic. Modbus TCP has no such watchdog: the HMI
simply retries on the next 500\,ms poll, the user sees a stale value for
one cycle, and nothing in the protocol forces a safe state.
Safety implication: a Profinet IO outage forces the process into its
engineered safe substitute state (valves close, drives coast to stop),
which is the desired behaviour on a safety-relevant interlock. A Modbus
outage hides the problem from the operator, who may act on stale data --
which is why Modbus must never be the path for a safety function.
\textit{Rubric: must mention the Profinet watchdog and the absence of one
in Modbus; deduct if the student claims Modbus also fails closed.}
\end{solutionbox}

\begin{solutionbox}
\textbf{Exercise 4.8 -- Compensating for Unauthenticated GOOSE.}
(1) \emph{Layer-2 segmentation}: put the GOOSE/SV bus on its own VLAN with
no routed access. Cheap, no protocol change. Drawback: a compromised
engineering laptop on the same VLAN bypasses it entirely.
(2) \emph{Bay-level monitoring}: deploy a passive IDS (e.g.\ a network tap
plus an OT IDS) that learns the expected GOOSE multicast MAC addresses,
state numbers, and sequence numbers, and alerts on duplicates or
unexpected publishers. Moderate cost, no latency impact. Drawback:
detection only -- the breaker has already tripped by the time the alert
arrives.
(3) \emph{IEC 62351-6 authentication}: enable HMAC-based message
authentication on every IED and merging unit. Highest cost (firmware
upgrades on every IED, key-management infrastructure, vendor support
matrix) and the protection engineer will object that the additional
processing pushes end-to-end latency past the 4\,ms budget for transfer
trip schemes.
\textit{Rubric: ordering may vary slightly, but the three controls must be
distinct (network, monitoring, crypto) and the drawback must be
operationally specific.}
\end{solutionbox}

\begin{solutionbox}
\textbf{Exercise 4.9 -- Spot the Bug in a Modbus Client.}
(1) The client opens and closes a fresh TCP connection every 100\,ms. On
the PLC side this consumes a slot in the (small) Modbus TCP connection
table -- typical S7-1200 or ControlLogix devices allow only 4 to 16
concurrent connections, and each closed connection sits in
\texttt{TIME\_WAIT} on the PLC stack for tens of seconds. Under sustained
polling the table fills, new connections are refused, and any other client
(an HMI, the engineering workstation) stops being able to talk to the PLC.
The operator sees the HMI go ``no data''.
(2) The standard fix is to open the TCP connection once and reuse it for
every transaction, multiplexing requests with the MBAP transaction ID. The
Modbus Organization explicitly recommends this in the
\emph{Modbus/TCP Implementation Guide} precisely because PLC connection
tables are small and \texttt{TIME\_WAIT} on embedded stacks is expensive.
\textit{Rubric: accept ``connection pooling'' as the fix; deduct if the
student blames bandwidth -- the bug is connection-table exhaustion, not
throughput.}
\end{solutionbox}

\begin{solutionbox}
\textbf{Exercise 4.10 -- Three Security Models Side by Side.}
\begin{center}\small
\begin{tabular}{p{2.4cm}p{3.6cm}p{3.6cm}p{3.6cm}}
\toprule
 & \textbf{OPC UA \texttt{Basic256Sha256}} & \textbf{MQTT over TLS} & \textbf{Profinet Security} \\
\midrule
Auth & X.509 mutual (application instance) $+$ user token & TLS server cert; optional mTLS; broker password & Certificate-based per IEC 62443, device identity \\
Integrity & RSA-SHA256 message signatures & TLS record MAC (AEAD) & HMAC over Profinet PDUs \\
Confidentiality & AES-256-CBC at message layer & TLS 1.2/1.3 AEAD (e.g.\ AES-GCM) & AES-GCM at PN frame layer \\
Key mgmt & PKI; optional GDS for cert push and CRL & Out-of-band PKI; broker-managed creds & Vendor PKI; engineering tool provisions certs \\
\bottomrule
\end{tabular}
\end{center}
OPC UA is hardest to operate at scale because each application instance
needs its own certificate, the trust lists must be kept consistent across
servers and clients, and revocation must propagate before a stolen cert
expires; without a GDS the work is manual per device. MQTT-TLS is the
easiest because the broker is a single choke point and standard web-PKI
tooling applies. Profinet Security sits in between but is constrained by
vendor support -- only newer IO devices implement it.
\textit{Rubric: cell entries may use different but equivalent mechanism
names; the operability paragraph must identify certificate lifecycle as the
OPC UA pain point.}
\end{solutionbox}

\end{document}
